\documentclass[10pt,a4paper]{amsart}
\usepackage[margin=19mm]{geometry}
\usepackage{amsmath,amssymb,mathtools}
\usepackage[scr=rsfs,cal=euler]{mathalfa}
\usepackage{xfrac}
\usepackage[unicode,hidelinks,bookmarks=false]{hyperref}
\newcommand{\ie}{i.e.\ }
\newcommand{\cf}{cf.\ }
\newcommand{\resp}{respectively\ }
\newcommand{\Subsection}{\subsection}
\providecommand{\nobreakdash}{\nobreak\hbox{-}}
\setlength{\emergencystretch}{2em}
\multlinegap0pt
\usepackage{amssymb}
\newtheorem{lemma}{Lemma}
\title{Complete Spectrum and Sharp Local Stability for the Critical Exponential Biharmonic Choquard Equation in \(\mathbb R^{4}\)}
\author{Wenjing Chen, Shengbing Deng}
\date{Source: arXiv:2608.01380v1 (CC BY 4.0)}
\begin{document}
\maketitle

\noindent\emph{Source and licence.} Complete Spectrum and Sharp Local Stability for the Critical Exponential Biharmonic Choquard Equation in \(\mathbb R^{4}\). Version arXiv:2608.01380v1 (CC BY 4.0), distributed by arXiv under the Creative Commons Attribution 4.0 International licence, http://creativecommons.org/licenses/by/4.0/, as recorded in the arXiv licence metadata for this exact version and shown on the abstract page https://arxiv.org/abs/2608.01380v1. Full licence notice: \texttt{licenses/arxiv-2608.01380-CC-BY-4.0.txt}.\par\smallskip

\noindent\textit{Editorial scope.} Counted unit: the article's lemma lem:normal-asymptotics (source0.tex lines 1271--1300). The article's complete original proof of it is delivered with it (source0.tex lines 1302--1479, one \texttt{proof} environment of 178 lines). No mathematics is written, repaired or re-derived: the statement and the proof are the article's own lines, in the article's own notation, and no retained line is substituted or re-flowed.  Retained uncounted as the source-local context the proof needs: lines 128--132; lines 140--142; lines 520--533.  Nothing was omitted from any retained span, so the package carries no omission record.  No retained line contains a citation, so the wrapper carries no bibliography and no bibliography entry.  No retained line contains a figure environment, an image-inclusion command or drawing code, so no image is delivered and the package declares no asset dependency.  Editorial formatting only, in addition: an AMS \texttt{amsart} preamble that restates the article's own declarations and macros used by the retained lines, namely the environments \texttt{lemma} on separate counters, together with the standard packages those lines need.  The retained environments print in this document as Lemma 1 in order of appearance, whereas the article prints the same items with the numbering of its own full document.  The complete native source of the article as distributed in the arXiv e-print for this exact version is bundled unchanged as \texttt{sources/206620/source0.tex}.
\begin{equation}\label{eq:main}
\Delta^{2}u=
\left(\int_{\mathbb R^{4}}\frac{e^{u(y)}}{|x-y|^{\alpha}}\,dy\right)e^{u(x)}
\qquad\hbox{in }\mathbb R^{4},
\end{equation}

\begin{equation}\label{eq:finite-mass}
\int_{\mathbb R^{4}}e^{\frac{8}{8-\alpha}u(x)}\,dx<+\infty.
\end{equation}

\begin{lemma}\label{lem:Riesz-estimate}
Let \(0<\alpha<4\), \(\theta>0\), and \(\theta+\alpha>4\). Then
\[
\int_{\mathbb R^{4}}
\frac{1}{|x-y|^{\alpha}}
\frac{1}{(1+|y|^{2})^{\theta/2}}\,dy
\leq C
\begin{cases}
(1+|x|^{2})^{(4-\alpha-\theta)/2},&\theta<4,\\
(1+|x|^{2})^{-\alpha/2}\log(2+|x|),&\theta=4,\\
(1+|x|^{2})^{-\alpha/2},&\theta>4.
\end{cases}
\]
\end{lemma}

\begin{lemma}\label{lem:normal-asymptotics}
Let \(u\) be a normal finite-mass solution of \eqref{eq:main}. Then
\begin{equation}\label{eq:strict-mass-lower}
\frac{p_{\alpha}M}{8\pi^{2}}>4,
\qquad p_{\alpha}:=\frac{8}{8-\alpha}.
\end{equation}
There exists \(\varepsilon_{1}>0\) such that
\begin{equation}\label{eq:F-tail-decay}
F(x)\leq C(1+|x|)^{-4-\varepsilon_{1}}
\qquad\hbox{for }|x|\geq1.
\end{equation}
Moreover, as \(|x|\to+\infty\),
\begin{equation}\label{eq:u-asymptotic}
u(x)=-\frac{M}{8\pi^{2}}\log|x|+O(1),
\end{equation}
\begin{equation}\label{eq:du-asymptotic}
\partial_{r}u(x)=-\frac{M}{8\pi^{2}}\frac1{|x|}+o(|x|^{-1}),
\end{equation}
\begin{equation}\label{eq:drr-u-asymptotic}
\partial_{rr}u(x)=\frac{M}{8\pi^{2}}\frac1{|x|^{2}}+o(|x|^{-2}),
\end{equation}
\begin{equation*}
\Delta u(x)=-\frac{M}{4\pi^{2}}\frac1{|x|^{2}}+o(|x|^{-2}),
\end{equation*}
and
\begin{equation}\label{eq:dr-lap-u-asymptotic}
\partial_{r}\Delta u(x)=\frac{M}{2\pi^{2}}\frac1{|x|^{3}}+o(|x|^{-3}).
\end{equation}
All the remainder estimates in \eqref{eq:du-asymptotic}--\eqref{eq:dr-lap-u-asymptotic} are uniform with respect to the angular variable \(x/|x|\).
\end{lemma}

\begin{proof}
By normality,
\begin{equation}\label{eq:normal-representation-asymptotic}
u(x)=-\frac{1}{8\pi^{2}}
\int_{\mathbb R^{4}}\log|x-y|F(y)\,dy+C.
\end{equation}
We first record the logarithmic moment. Fix \(x_{0}\in\mathbb R^{4}\). On a neighbourhood of \(x_{0}\), the negative part of \(\log|x_{0}-y|\) is integrable against \(F\), because \(F\in L^{q}_{\operatorname{loc}}\) for every finite \(q>1\). For \(|y|\) large,
\(\log|x_{0}-y|\geq\log(2+|y|)-C\). Since the integral in \eqref{eq:normal-representation-asymptotic} is finite, it follows that
\begin{equation}\label{eq:logarithmic-moment}
\int_{\mathbb R^{4}}\log(2+|y|)F(y)\,dy<+\infty.
\end{equation}
The elementary bound
\[
|x-y|\leq(1+|x|)(1+|y|)
\]
and \eqref{eq:logarithmic-moment} give
\begin{equation}\label{eq:lower-u-preliminary}
u(x)\geq-\gamma\log(1+|x|)-C.
\end{equation}
If \(p_{\alpha}\gamma\leq4\), then \eqref{eq:lower-u-preliminary} would imply
\(e^{p_{\alpha}u(x)}\geq C(1+|x|)^{-p_{\alpha}\gamma}\), which is not integrable in \(\mathbb R^{4}\). This contradicts \eqref{eq:finite-mass} and proves \eqref{eq:strict-mass-lower}.

We next obtain an almost sharp upper bound without assuming uniform local control of \(F\) at infinity. Fix
\begin{equation}\label{eq:epsilon-choice-asymptotic}
0<\varepsilon<\gamma-\frac{8-\alpha}{2}.
\end{equation}
Choose \(s>\max\{1,4/(4-\alpha)\}\). Since \(F\in L^{1}(\mathbb R^{4})\), one can choose \(R>1\) such that, with
\[
m_{R}:=\int_{\mathbb R^{4}\setminus B_{R}}F(y)\,dy,
\qquad
\gamma_{R}:=\frac{1}{8\pi^{2}}\int_{B_{R}}F(y)\,dy,
\]
we have
\begin{equation}\label{eq:tail-small-choice}
\gamma_{R}>\gamma-\varepsilon,
\qquad
\frac{s m_{R}}{8\pi^{2}}<4.
\end{equation}
Let \(|x|>2R+8\). Define
\[
v_{x}(z):=-\frac{1}{8\pi^{2}}
\int_{B_{4}(x)}\log|z-y|F(y)\,dy,
\qquad
h_{x}:=u-v_{x}.
\]
Then \(\Delta^{2}h_{x}=0\) in \(B_{4}(x)\). For \(z\in B_{2}(x)\), the set \(B_{R}\) is disjoint from \(B_{4}(x)\), and
\[
\log|z-y|=\log|x|+O_{R}(1)
\qquad\hbox{uniformly for }y\in B_{R}.
\]
If \(y\notin B_{4}(x)\), then \(|z-y|\geq2\). The contribution of
\(\mathbb R^{4}\setminus(B_{R}\cup B_{4}(x))\) to \(h_x(z)\) is
nonpositive. The fixed additive constant in \eqref{eq:normal-representation-asymptotic}
and the bounded error on \(B_R\) can be absorbed into \(C_R\). Hence
\begin{equation}\label{eq:hx-upper}
h_{x}(z)\leq-\gamma_{R}\log|x|+C_{R}
\qquad\hbox{for }z\in B_{2}(x).
\end{equation}

Put \(m_{x}:=\int_{B_{4}(x)}F\leq m_{R}\). If \(m_x>0\), let
\(d\nu_x=F(y)m_x^{-1}dy\) on \(B_4(x)\). For \(z\in B_2(x)\), Jensen's inequality yields
\[
\begin{aligned}
e^{s v_x(z)}
&\leq
\exp\left(\frac{s m_x}{8\pi^{2}}
\int_{B_4(x)}\log\frac{6}{|z-y|}\,d\nu_x(y)\right)\leq
\int_{B_4(x)}
\left(\frac{6}{|z-y|}\right)^{s m_x/(8\pi^{2})}
\,d\nu_x(y).
\end{aligned}
\]
The same conclusion is immediate if \(m_x=0\). Integrating in \(z\), using \eqref{eq:tail-small-choice}, and then using \eqref{eq:hx-upper}, we obtain
\begin{equation}\label{eq:uniform-moving-exponential}
\int_{B_{2}(x)}e^{s u(z)}\,dz
\leq C_{R}|x|^{-s\gamma_{R}}.
\end{equation}

For \(z\in B_{1}(x)\), split the Riesz potential as
\[
V(z)=\int_{B_{2}(x)}\frac{e^{u(y)}}{|z-y|^{\alpha}}\,dy
+\int_{\mathbb R^{4}\setminus B_{2}(x)}
\frac{e^{u(y)}}{|z-y|^{\alpha}}\,dy.
\]
The first term is bounded by \(C_R|x|^{-\gamma_R}\) by
\eqref{eq:uniform-moving-exponential} and the choice of \(s\). For the second term, let \(p_{\alpha}'=8/\alpha\). Since \(|z-y|\geq1\) there and \(\alpha p_{\alpha}'=8>4\), H\"older's inequality and \eqref{eq:finite-mass} give a bound independent of \(x\). Thus
\begin{equation}\label{eq:V-moving-bound}
\|V\|_{L^{\infty}(B_{1}(x))}\leq C_{R}.
\end{equation}
Combining \eqref{eq:uniform-moving-exponential} and \eqref{eq:V-moving-bound},
\begin{equation*}
\|F\|_{L^{s}(B_{1}(x))}\leq C_{R}|x|^{-\gamma_{R}}.
\end{equation*}
Moreover, \(v_x\geq-m_R\log 6/(8\pi^{2})\), while the preceding Jensen estimate controls its positive part. Hence \(\|v_x\|_{L^{s}(B_1(x))}\leq C_R\). The interior \(W^{4,s}\)-estimate for
\(\Delta^{2}v_x=F\) in \(B_1(x)\) now gives
\[
\|v_x\|_{L^{\infty}(B_{1/2}(x))}
\leq C\left(\|v_x\|_{L^{s}(B_1(x))}
+\|F\|_{L^{s}(B_1(x))}\right)\leq C_R.
\]
Together with \eqref{eq:hx-upper}, this proves
\begin{equation*}
u(x)\leq-(\gamma-\varepsilon)\log|x|+C_{\varepsilon}
\qquad\hbox{for }|x|\geq2.
\end{equation*}

Choose \(\varepsilon\) in \eqref{eq:epsilon-choice-asymptotic} and set
\(\theta:=\gamma-\varepsilon>(8-\alpha)/2\). Then
\[
e^{u(x)}\leq C(1+|x|^{2})^{-\theta/2}.
\]
Since \(\theta+\alpha>4\), Lemma~\ref{lem:Riesz-estimate} applies. If \(\theta<4\), it gives
\[
F(x)\leq C(1+|x|^{2})^{(4-\alpha-2\theta)/2}.
\]
If \(\theta=4\), the same estimate holds with an additional logarithm; if \(\theta>4\), it gives
\(F(x)\leq C(1+|x|^{2})^{-(\alpha+\theta)/2}\). In the first case the corresponding power of \(|x|\) is
\(4-\alpha-2\theta<-4\). In the borderline case the logarithm is
absorbed by a smaller negative power. In the last case the power is
\(-\alpha-\theta<-4\). Thus, after decreasing a positive exponent if
necessary, \eqref{eq:F-tail-decay} follows.

We can now return to the logarithmic potential without a moving singularity problem. Let \(R=|x|\) and \(e=x/R\). Split
\[
\int_{\mathbb R^{4}}\log\frac{|x-y|}{R}F(y)\,dy
=I_{1}(x)+I_{2}(x),
\]
where \(I_1\) is over \(|y|\leq R/2\). On this region
\(|\log|e-y/R||\leq C\), so \(|I_1(x)|\leq CM\). By \eqref{eq:F-tail-decay} and the change of variables \(y=Rz\),
\[
|I_2(x)|
\leq C R^{-\varepsilon_{1}}
\int_{|z|>1/2}|\log|e-z||\,|z|^{-4-\varepsilon_{1}}\,dz
\leq C R^{-\varepsilon_{1}}.
\]
The last integral is finite uniformly for \(e\in\mathbb S^{3}\). Formula
\eqref{eq:normal-representation-asymptotic} therefore proves
\eqref{eq:u-asymptotic}.

Finally, \eqref{eq:F-tail-decay} and local smoothness imply that the first three derivatives of the logarithmic potential may be taken under the integral sign. For \(m=1,2,3\),
\begin{equation}\label{eq:scaled-tail-kernels}
R^{m}\int_{|y|>R/2}|x-y|^{-m}F(y)\,dy
\leq C R^{-\varepsilon_{1}}
\int_{|z|>1/2}|e-z|^{-m}|z|^{-4-\varepsilon_{1}}\,dz
=o(1).
\end{equation}
On \(|y|\leq R/2\), the relevant scaled kernels are uniformly bounded because
\(|e-y/R|\geq1/2\). For each fixed \(A>0\), their limits on \(|y|\leq A\)
are uniform in \(e=x/R\in\mathbb S^{3}\):
\[
R\frac{(x-y)\cdot x}{|x-y|^{2}|x|}\longrightarrow1,
\]
\[
R^{2}\partial_{rr}\log|x-y|
=R^{2}\left(
\frac{1}{|x-y|^{2}}
-2\frac{((x-y)\cdot e)^{2}}{|x-y|^{4}}
\right)\longrightarrow-1,
\]
\[
R^{2}\Delta\log|x-y|\longrightarrow2,
\qquad
R^{3}\partial_{r}\Delta\log|x-y|
=-4R^{3}\frac{(x-y)\cdot x}{|x|\,|x-y|^{4}}
\longrightarrow-4.
\]
On the remaining region \(A<|y|\leq R/2\), the uniform boundedness of these
kernels gives an error bounded by
\[
C\int_{|y|>A}F(y)\,dy.
\]
We first let \(R\to+\infty\) with \(A\) fixed and then let \(A\to+\infty\).
This proves uniform convergence of all interior contributions. The tail
estimate \eqref{eq:scaled-tail-kernels} applies to the complementary region.
Combining these facts with \eqref{eq:normal-representation-asymptotic} yields
\eqref{eq:du-asymptotic}--\eqref{eq:dr-lap-u-asymptotic}, including
\eqref{eq:drr-u-asymptotic}, with uniform remainders in the angular variable.
\end{proof}

\end{document}
